\subsection{Data integrity and causal trace construction}
\label{sec:C07:data}
Deduplication is a statistical control as well as a storage optimization. Repeated content changes effective sampling weights and can leak benchmark material across splits. Exact hashing catches identical bytes; normalized or approximate matching catches additional overlap but can merge legitimately distinct examples. For code and agent traces, repository lineage, issue templates, and generated variants can carry the same underlying solution even when strings differ. The deduplication policy should state its unit and false-positive tradeoff. A claim of uncontaminated evaluation requires evidence about the held-out tasks, not merely that the training corpus passed a generic duplicate filter.

Agent trace construction needs stronger provenance than ordinary document collection. Each observation should identify its source and whether it was actually available before the action. A later grader's explanation is not an environment observation. A synthetic tool result is not a measured one. A trace modified after execution can be useful training data, but its conditioning semantics must be reconstructed deliberately. MemoryWalker directly studies training on contexts that differ from the live contexts of edited agent rollouts~\cite{j2026memorywalker}. Its relevance extends to pretraining mixtures containing retrospective summaries or reconstructed tool histories, even when no reinforcement-learning update is being performed.

One practical audit reconstructs the visible prefix for every supervised assistant action and compares it with the serialized training input. Differences are classified as intentional transformations, unavailable evidence, or unresolved discrepancies. The audit should also check whether system instructions, tool schemas, and permissions match the recorded environment. A policy trained to imitate an action that was valid under an earlier schema can emit an invalid operation after deployment changes. Data age alone does not identify this mismatch; compatibility metadata does.

Curriculum and filtering can introduce selection effects. Retaining only successful trajectories teaches useful behavior but removes evidence about failure recognition and recovery. Retaining only traces accepted by one learned judge can amplify that judge's blind spots. A mixture containing both successful and failed traces needs explicit targets: imitation of an erroneous action is not made desirable merely because the surrounding record is informative. Alternative objectives can predict outcomes, classify failure states, or train corrective continuations. The experiment should specify which information is learned from each trace category rather than relying on a vague label such as high-quality data.

